\documentclass[11pt]{article}
\usepackage[margin=1in]{geometry}
\usepackage{tabularx}
\usepackage[hidelinks]{hyperref}
\setlength{\emergencystretch}{2em}
\title{Student Roster Form Automation\\\large Project Proposal and Development Report}
\author{Michael Kowata\\ITC-492 Senior Capstone}
\date{Draft updated October 28, 2026}
\begin{document}
\begin{titlepage}
\maketitle
\thispagestyle{empty}
\begin{abstract}
This project proposes a Google Apps Script application that converts a screenshot
of Zoom's Participants panel into a reviewed student roster and a separate Google
Form for each approved roster entry. An instructor will upload the screenshot
through a sidebar in a control spreadsheet, correct extracted names, and approve
generation of the output spreadsheet and personalized form copies. The project
targets repetitive transcription, copying, and link-matching work while retaining
human review where optical character recognition (OCR) may be inaccurate.
This report documents the accepted baseline, planned architecture, requirements,
and evaluation method. The proposal presentation has been reported as successful
by the project owner. Application implementation and measured evaluation results
are not yet present in the repository reviewed for this draft.
\end{abstract}
\vfill
\noindent\textbf{Document status:} Living proposal and development report.
Planned behavior is distinguished from completed work throughout the paper.
\end{titlepage}
\tableofcontents
\clearpage
\section{Introduction}
\subsection{Problem and Motivation}
An instructor who needs an individual form for every student must transcribe
names, duplicate a form template, and place the correct responder link beside
each student. Repeating these steps creates opportunities for omitted entries,
duplicate copies, and mismatched links. The proposed application will automate
the repeated operations while allowing the instructor to verify the roster.
Reduced effort is an intended benefit, not a measured result of this proposal.

\subsection{Objective and Intended User}
The intended user is an instructor operating one Google account with access to
a form template. The objective is a complete workflow from one supported image
to an output roster in which every approved entry has its own form and matching
responder link. A shared form with prefilled names does not satisfy this goal:
the project requires distinct forms and independent response sets.

\subsection{Scope}
The initial input is one PNG or JPEG screenshot of Zoom's Participants panel.
Only visible names are included. The operator can correct display names and
remove non-students; the application will not infer identities from nicknames.
The first version supports one operator, one account, one template, and one
roster per run. Gallery-view tiles, scrolling, combining screenshots, automatic
screen capture, handwriting, email delivery, grading, and response aggregation
are outside the initial scope. These boundaries keep the capstone focused on
one testable task rather than a general classroom-management platform.

\section{Technical Background and Related Work}
\subsection{Scripting Platform}
Google Apps Script supports JavaScript through its V8 runtime. It is not a
standard Node.js or browser runtime, so libraries must be evaluated for that
environment rather than assumed compatible~\cite{googleAppsScriptV8}.
This project therefore adopts Apps Script instead of the earlier Python
command-line assumption.

Google's HTML service supports asynchronous server calls through
\texttt{google.script.run}. A form containing a file-input field can pass the
selected file to a server function as a Blob~\cite{googleHtmlCommunication}.
This supports the proposed local-file upload without requiring the instructor
to manually upload an image to Drive first. The chosen OCR method may still
need temporary cloud files; upload convenience does not establish a
storage-free processing pipeline.

These platform references are technical documentation, not experimental research
on OCR accuracy. The following research provides context for selecting and
evaluating an OCR approach; it does not establish performance on Zoom screenshots.

\subsection{Traditional OCR Pipelines}
Smith's conference paper describes the historical Tesseract engine, including
line finding, word recognition, and adaptive classification~\cite{smith2007tesseract}.
It illustrates that OCR involves locating and organizing text as well as
recognizing characters. The paper concerns the engine described in 2007 and
should not be treated as a specification or benchmark of current Tesseract releases.

For this capstone, the design implication is to evaluate the complete extraction
pipeline. A recognizer may read interface text correctly but still produce an
incorrect student roster if controls or role labels are included. This is our
application-specific inference, not a result reported by Smith. Trials should
therefore distinguish missed name regions, character errors, and unwanted
interface text, and compare an uncropped panel with a name-region crop.

\subsection{Neural Sequence Recognition}
Shi, Bai, and Yao propose an architecture combining feature extraction, sequence
modeling, and transcription for scene-text recognition~\cite{shi2015crnn}.
Their approach handles variable-length sequences without explicit character
segmentation and supports recognition without a predefined lexicon. The cited
version is the authors' 2015 arXiv preprint; its evaluations concern scene-text
benchmarks rather than Zoom participant lists.

This motivates testing names of different lengths and spellings instead of
assuming a fixed dictionary will cover a roster. It does not justify automatically
replacing an unfamiliar name with a common word. Human review remains necessary.
Training a new recognition model is outside this project's initial scope; this
research informs evaluation of existing tools.

\subsection{Practical OCR Systems}
Du et al. describe PP-OCR, a lightweight system with text detection, direction
classification, and recognition components~\cite{du2020ppocr}. Their work studies
model capability and size through optimization strategies and ablation experiments.
The cited source is a 2020 arXiv preprint and describes that version of PP-OCR,
not every later PaddleOCR release.

The relevant lesson for this project is to consider processing cost alongside
recognition quality. We will record time to produce a reviewable roster and
manual correction effort, rather than select a tool solely by an unrelated
benchmark score. A Zoom panel has a narrower layout than the broader text-image
settings addressed by this research; local experiments remain necessary.

\subsection{Connecting OCR to JavaScript Automation}
Tesseract.js provides a JavaScript interface to a WebAssembly port of Tesseract
for browsers and Node.js~\cite{tesseractJs}. This is primary software documentation,
not a peer-reviewed accuracy study. It offers a candidate scripting route:
test recognition in the browser-based sidebar and pass extracted text to the
Apps Script workflow. Compatibility, model loading, worker behavior, and latency
must be demonstrated in the actual sidebar before adopting this approach.

Browser or Node.js support does not imply compatibility with server-side
\texttt{Code.gs}; Apps Script's runtime differs from both~\cite{googleAppsScriptV8}.
An alternative architecture is to upload the image for server-orchestrated OCR.
M2 will compare feasible integration paths, including permissions, storage,
cleanup, and any service cost. No engine or deployment path is selected by this
literature review.

\subsection{Research Questions for the Technology Trial}
\begin{enumerate}
\item Does cropping to the name region reduce unwanted interface text without
losing student entries?
\item How do readable screenshot scale, name length, and duplicate display names
affect extraction errors and human correction effort?
\item Which feasible scripting integration produces a correctable roster with
acceptable latency and manageable authorization and data handling?
\end{enumerate}
These questions translate the literature into project-specific experiments.
They are proposed comparisons, not claims of published or measured improvements.

\clearpage
\section{Requirements and Data Handling}
\subsection{Functional Baseline}
The requirements below summarize the identifiers maintained in
\texttt{docs/requirements.md}. Acceptance is based on observable behavior.

\begin{center}
\renewcommand{\arraystretch}{1.2}
\begin{tabularx}{\textwidth}{l X}
\hline
\textbf{IDs} & \textbf{Required behavior} \\
\hline
R01--R02 & Validate a local sidebar upload, extract visible names, and allow
correction, removal, and addition before explicit approval. \\
R03--R04 & Preserve a stable identity for every entry, including identical
display names, and generate one output row per approved entry. \\
R05--R07 & Create a distinct template copy for every entry, preserve questions
and independent responses, and write the matching responder URL. \\
R08--R09 & Resume interrupted work without extra copies and retain actionable
status information for partially completed entries. \\
R10 & Protect private artifacts and document cloud access and temporary-resource
cleanup. Use fictional names in committed test fixtures. \\
\hline
\end{tabularx}
\end{center}

\subsection{Approval and Privacy}
The control spreadsheet is the working surface; the generated roster is a
separate deliverable. Cancellation before approval must not create an output
roster or student forms. Temporary OCR resources may already exist and must
follow a documented cleanup policy on success, failure, and cancellation.

Real screenshots, rosters, credentials, and exported run state must not enter
version control. Local ignore rules support this policy but do not replace
review of staged files. Cloud permissions, retention, and cleanup mechanisms
remain design decisions to validate. A form bearing a student's name does not
authenticate the respondent; access settings must be tested separately.

\section{Proposed Design}
\subsection{User Workflow}
\begin{enumerate}
\item Open a custom menu in the control spreadsheet and display the sidebar.
\item Select a local Zoom Participants screenshot and validate the upload.
\item Extract visible names, excluding interface controls and role/self labels.
\item Review and correct the roster in a staging tab, then explicitly approve it.
\item Create the separate roster spreadsheet and copy the template for each entry.
\item Write each responder link and completion status to its corresponding row.
\end{enumerate}

\subsection{Components and Recovery}
Planned source files include \texttt{Code.gs} for server operations,
\texttt{Sidebar.html} for the upload interface, and \texttt{appsscript.json}
for configuration. These are planned components, not implemented modules.
Persistent run records should associate a run identifier and stable entry
identifier with the name, form identifier, responder URL, and processing status.
Names alone cannot be keys because two students may share the same display name.

A failed request after copying a form but before recording its identifier can
cause a duplicate on retry. The implementation must reconcile existing resources
before repeating creation; a local success flag alone is insufficient. The
storage mechanism and reconciliation method will be selected and tested in M2.
Upload limits, authorization scopes, execution limits, and OCR selection are
also open technical decisions.

\clearpage
\section{Evaluation Method}
\subsection{Fixtures and Measurements}
Use a Participants panel sample containing fictional names, including duplicate
display names. Target 20 visible entries only if they remain legible in one
screenshot; otherwise record the actual visible count. Maintain an independently
checked expected roster. Compare raw OCR with that roster before correction,
reporting exact-name matches, missing and extra entries, and the number of manual
corrections. No accuracy threshold is claimed before an initial OCR trial.

For the comparisons motivated by the literature, keep the expected roster fixed
and vary one input condition at a time, such as cropping or readable image scale.
Preserve duplicate entries when matching outputs to expected rows. Record engine
and model versions, language settings, preprocessing, and browser or service
configuration. Separate initial model-loading time from subsequent recognition
time when relevant. Report each tested condition and failure rather than
generalizing from the best sample to all participant screenshots.

After review, let $N$ denote the number of approved entries. Require exactly
$N$ distinct forms and $N$ correct name-to-link mappings. Check copied questions
against the template and submit fictional responses to verify that each response
appears only in its corresponding form. Test responder access using a
representative intended respondent rather than relying on the owner's editor view.

\subsection{Failure Scenarios and Reproducibility}
Test missing, invalid, empty, and oversized uploads; cancellation; duplicate
names; and partial service failures. Interrupt runs after spreadsheet creation,
form copying, state saving, and row updates. Resuming the same run must complete
remaining work without creating additional forms. Record cleanup outcomes for
temporary OCR resources in both successful and failed runs.

For each evaluation, record the source commit, account configuration, template
settings, fixture identifier, approved entry count, commands or UI steps, and
actual outputs. Measure total runtime and report limitations. If time savings
are evaluated, compare manual and automated workflows on the same fixture,
including human correction time; do not infer savings from successful execution.

\section{Results and Discussion}
\subsection{Current Evidence}
The repository contains a defined scope, acceptance criteria, GitHub issue and
milestone tracking, published Wiki pages, experiment templates, an agent log,
and presentation artifacts. These demonstrate project preparation, not the
correctness of the proposed automation. This draft contains no measured OCR
accuracy, runtime, successful form-generation count, or recovery results.

\subsection{Presentation Feedback}
The proposal presentation was scheduled for September 28, 2026. On October 28,
the project owner reported that it went successfully. The reported feedback was
to write the paper as the project develops and to organize it as a report across
multiple pages rather than leave it as a single-page outline. This revision
responds by expanding the proposal, separating major topics, and establishing
an ongoing documentation process. The actual delivery date and demonstrated
commit have not been independently confirmed in this draft.

\subsection{Limitations}
Image quality, display-name ambiguity, service permissions, and execution limits
may constrain the solution. Human review reduces the effect of extraction errors
but introduces work that must be included in evaluation. A successful sample
does not establish support for every Zoom layout or account configuration.

\clearpage
\section{Development and Documentation Plan}
M1 establishes scope and the project foundation. M2 evaluates upload, OCR,
authorization, and storage choices. M3 implements extraction and review; M4
creates the roster and individual forms; M5 addresses links and recovery.
M6 evaluates and documents the system, and M7 prepares final delivery. Milestone
dates are planning targets, not evidence that these activities are complete.

Update this paper alongside each implementation or experiment change. For a
design decision, record the alternatives, rationale, and associated issue. For
an experiment, add its method and actual evidence before drawing conclusions.
For an implemented component, replace proposed behavior with a description of
the verified implementation and its source commit. Add verified bibliography
entries when sources are consulted, and retain unsuccessful results and known
limitations rather than reporting only successful examples.

\section{Conclusion and Future Work}
The proposal defines a bounded automation task: transform visible participant
names into a reviewed roster and independent student forms with matching links.
Its main design commitments are an Apps Script sidebar upload, explicit human
approval, stable entry identities, and recoverable creation of output resources.
The next technical work is to establish the feasibility of screenshot extraction
and Google integration through documented trials. Conclusions about accuracy,
time savings, and reliability must wait for implementation and measurement.

\clearpage
\bibliographystyle{plain}
\bibliography{references}
\clearpage
\appendix
\section{Supporting Artifacts}
\begin{itemize}
\item Repository: \url{https://github.com/Michael-Kowata/ITC-492}
\item Wiki: \url{https://github.com/Michael-Kowata/ITC-492/wiki}
\item Requirements: \texttt{docs/requirements.md}; proposal: \texttt{docs/proposal.md}.
\item Experiment reports: \texttt{sandbox/}; no completed experiment is asserted here.
\item Proposal slides and notes: \texttt{presentations/2026-09-28-proposal/}.
\item Assistance and review evidence: \texttt{AGENT-LOG.md}.
\end{itemize}
The final report should identify the evaluated release and link sanitized raw
results. Until then, repository paths identify evolving artifacts rather than
a completed software release.
\end{document}
